\documentclass[a4paper]{article}

\usepackage{mlmodern}
\usepackage[utf8]{inputenc}
\usepackage[T1]{fontenc}
\usepackage{textcomp}
\usepackage{amsmath, amssymb}
\edef\restoreparindent{\parindent=\the\parindent\relax}
\DeclareSymbolFont{letters}{OML}{ztmcm}{m}{it}
\usepackage{parskip}
\restoreparindent
\newtheorem{problem}{Problem}
\newtheorem{example}{Example}
\newtheorem{lemma}{Lemma}
\newtheorem{theorem}{Theorem}
\newtheorem{problem}{Problem}
\newtheorem{example}{Example}
\newtheorem{definition}{Definition}
\newtheorem{lemma}{Lemma}
\newtheorem{theorem}{Theorem}

\begin{document}
    

Hello Jen,

I hope you'll excuse the unavoidable corniness that comes with written letters.
Though the future is always uncertain, it is conceivable that Baltimore was the
last place we'll see each other in person for quite some time. With my time at
the lab nearing its end, I feel obliged to express to you what your mentorship
and the work we've done together has meant to me. While this is mostly inspired
by my own desire to make my gratitude known, I also want to leave no doubt in
your mind about how extraordinary this experience has been for me.

You do not know this, as the boundaries of our professional relationship gave
no opportunity to share it, but the years prior to entering the lab were
extraordinarily difficult for me. I was navigating a personal tragedy involving
the passing of a partner as well as heavy family caretaking responsibilities right
as the pandemic was at its peak. So I did not have the expected early twenties experience,
so to speak.

It isn't necessary to dwell on how things slowly improved from there. Suffice it
to say that entering your lab was, to me, like the final stone needed to rebuild
the fallen structures. I had dreamed for so long of working in science, living
with the contradiction of believing in my own potential while feeling that life
wasn't giving me an opportunity, and suddenly I was finally given a chance! When
I was going through those dark times, I would have never, not in a million
years, imagined that a few years later I'd be working at UPenn, publishing
research, traveling to conferences, and meeting such incredible people. I would
have laughed at anyone who suggested such a thing!

I know you are very aware of the academic impact your mentorship and the
opportunities you've given me have had. But I don't think you know how much it
meant for my life in general. Without knowing me, and without needing to, you
gave me a chance, precisely the thing I needed the most. I've fallen in love
with sleep science ever since, something I did not even know existed prior to
knowing you, and it feels so good to finally have something I can call «my
thing».

I want to thank you again for everything you've done for me, from this last trip
to Philly and Baltimore all the way back to that initial email where you decided
to give this random Argentinian guy an opportunity. I regret very much not
realizing that the day of our presentations was the last day we'd see each other
in person. I owed you a more proper goodbye and a more proper thank you, so I
hope this letter can serve both purposes.

I hope you know that, no matter how many years go by, or where we are at any
point in life, you will always be my mentor.

~

\newline
\newline
---Santi






















\end{document}
